% 【基础正文】所属：第5章（由 chapters/revised/05-social-knowledge.tex \input 引入）
\minitopic{从个人主义到社会认识论}个体信念依赖语言、档案、仪器、教育和他人劳动。承认依赖不等于否定个人理性，而是改变分析单位。\term{社会认识论}{social epistemology}既研究个人在社会关系中的知识，也研究制度和群体的认识表现\parencite{goldman1999,longino1990}。 评价制度可以采用真理促进、错误发现、观点包容或认识公正等标准。同行评议可能筛除错误，也可能复制偏见；开放讨论可能汇集信息，也可能放大噪声。制度没有因“社会性”而自动不客观，客观性有时正来自公开批评和角色分工。

\minitopic{群体能否知道}一种聚合观把群体信念还原为成员信念，例如多数人相信$p$。但委员会可能按规则接受一个没有多数成员私下相信的结论；多数成员也可能各自相信$p$，群体却从未形成相关决定。\term{联合承诺}{joint commitment}或功能结构理论把群体态度连接到决策程序、信息流和行动能力。 群体知识需要至少三类条件：结论为真；群体有形成和维持该结论的适当机制；该结论能指导群体层面的说明、决策或修正。把公司拥有数据库就算作知道过于宽松，把每个成员必须知道全部内容又过于狭窄。

\minitopic{权力与认识不正义}弗里克区分两类\term{认识不正义}{epistemic injustice}\parencite{fricker2007}。\term{证言不正义}{testimonial injustice}发生在身份偏见使听者不当地降低说话者可信度时；\term{诠释不正义}{hermeneutical injustice}发生在共同解释资源存在结构性缺口，使某些经验难以被理解和表达。这里要注意概念范围：在弗里克的较窄框架中，并非每一次错误的可信度判断都自动构成证言不正义，身份偏见及其权力结构是关键；后续社会认识论则把讨论扩展到沉默、可信度过剩、制度性排除等更广机制。使用“认识不正义”时应说明采用的是原始分类还是扩展意义。 前者不只是礼貌问题：可信度分配错误会阻断知识流。后者也不只是缺少词汇，而与谁有机会参与概念生产有关。Dotson进一步讨论沉默实践：听者的无能或不愿倾听可使说话行为无法完成\parencite{dotson2011}。 认识公正不要求把身份本身当作主张为真的证据；它要求校正与相关能力、证据无关的身份偏差，改善参与条件，并使质疑标准对称适用。

\minitopic{多样性何时有认识价值}观点多样性可以揭示隐藏假设、扩大问题选择并提高错误检测，但只有在成员能实际发言、批评被记录、资源不被垄断时才发挥作用。象征性代表并不自动改变信息结构。反过来，多样性也不意味着每个观点权重相等；方法质量和相关专长仍重要。

\begin{argumentbox}
制度评价应避免两种化约：只问结论有多少为真，会忽略不公正取得真相的代价；只问参与是否平等，会忽略制度是否稳定地产生错误。成熟的社会认识论同时评估真理表现、纠错能力和可信度分配。
\end{argumentbox}

\begin{comparison}
儒家知识传统强调师友、礼制与共同学习，说明认识从不完全孤立；同时，角色秩序可能限制谁能发言。用社会认识论分析它，既不能把关系性浪漫化，也不能以现代个人主义简单否定。应具体考察发言权、纠错渠道和权威可问责性。
\end{comparison}

\minitopic{制度客观性}客观性不必理解为研究者没有立场。制度可以通过公开数据、可复现方法、对抗性批评和利益披露，使个体偏见更容易被发现\parencite{longino1990}。个人完全中立也许不可得，程序性的可批评性却可以改善共同结果。 指标会改变行为。只奖励论文数量可能诱发切分发表，只奖励显著结果会放大选择性报告。认识制度的设计需要预测激励反馈，而不能把规则当作被动容器。纠错速度、负面结果可见性和复核资源都是评价维度。

\minitopic{群体失败的类型}信息级联发生在后来者忽略私人证据、追随早期公开选择时；群体极化使讨论后立场更极端；共同知识缺失则可能让每个人都知道风险，却不知道其他人也知道，因而无人行动。三种失败机制不同，对策也不同：独立先行判断、匿名汇总、明确责任和公开不确定性分别针对不同结构。 “集体愚蠢”不应被解释成成员都无知。聪明成员在错误激励和沟通网络中也会产生糟糕结果。反过来，良好制度能让没有任何全知成员的群体获得可靠结论。

\minitopic{案例工作坊：事故调查}工程师各自发现一个异常，却分别认为问题属于别的部门。管理层只接收汇总后的绿色指标，最终系统失效。个体层面没有人掌握完整风险，组织层面也没有机制把分散信息组合成警报。事故后不能只寻找一个“没看见真相的人”，而应审查上报通道、指标压缩和提出坏消息的代价。 改进包括记录异议、为跨部门异常设立汇合点、保护提出反证者，并让决策者说明忽略警报的理由。这些制度既分配责任，也构成群体获得知识的认识机制。

\minitopic{群体主体主张有解释收益也有边界}把知识归给委员会或机构，可以解释无人单独掌握全部材料时，组织为何仍能回答问题、作出承诺并承担修正责任。但“群体知道”不能遮蔽内部差异：某部门可能握有反证却无法上报，正式记录也可能与实际行动脱节。每次使用群体归属，都应说明知识由何种程序整合、在哪个接口可被调用，以及哪些成员缺乏访问。若这些问题没有答案，群体语言可能只是把组织名称放在无人负责的位置。
